\section{Conclusion}

We replicate a fixed equity volatility-routing method in a broad Chinese futures
cross-section. The primary market proxy produces a superficial echo of the source result:
the router has a slightly higher endpoint than momentum, and one high-volatility event is
successfully hedged by reversal. The broader evidence rejects a stable replication. The
increment is statistically negligible, reversal wins in a minority of Q4 months, the
representative Q4 observation favors momentum, and an equal-product volatility proxy
reverses the key conditional ranking.

At the same time, high volatility is economically meaningful. It identifies months with
higher cross-product correlation, dispersion, and directional coherence. The failure is
therefore not that volatility measures nothing. It is that shock intensity is not a
sufficient statistic for shock type.

This distinction explains why a scalar state can be effective in one asset class and fail
in another without contradiction. In an equity market, high volatility may reliably
proxy a distress, funding, liquidity, and rebound configuration. In a heterogeneous
futures market, high volatility can arise from physical scarcity, policy repricing,
macro risk, or forced trading, with different implications for continuation and reversal.

The appropriate conclusion is narrow: the published volatility-only router does not
robustly transport to this Chinese-futures sample under the pre-specified mapping. The
result does not justify a new optimized percentile. It motivates a more disciplined
research design in which direction, breadth, and stage are treated as explicit hypotheses
and any resulting router is evaluated only after the mechanism is frozen.
